\chapter{1997年全国卷（理）}
\section{单选题}

\begin{problem}
设集合 \(M = \{x \mid 0 \leqslant x < 2\}\)，集合 \(N = \{x \mid x^2 - 2x - 3 < 0\}\)，集合 \(M \cap N =\)（\(\quad\)）

\choices
    {\(\{x \mid 0\leqslant x < 1\}\)}
    {\(\{x \mid 0\leqslant x < 2\}\)}
    {\(\{x \mid 0\leqslant x\leqslant 1\}\)}
    {\(\{x \mid 0\leqslant x\leqslant 2\}\)}
\end{problem}

\begin{answer}
B
\end{answer}

\begin{solution}
由 \(x^{2}-2x-3=(x-3)(x+1)\)，得 \(N=(-1,3)\). 因而 \(M\cap N=[0,2)\)，故选 B.
\end{solution}

\begin{problem}
如果直线 \(ax + 2y + 2 = 0\) 与直线 \(3x - y - 2 = 0\) 平行，那么系数 \(a =\)（\(\quad\)）

\choices
    {\(-3\)}
    {\(-6\)}
    {\(-\frac{3}{2}\)}
    {\(\frac{2}{3}\)}
\end{problem}

\begin{answer}
B
\end{answer}

\begin{solution}
两直线分别化为 \(y=-\dfrac{a}{2}x-1\) 和 \(y=3x-2\)，它们平行，所以斜率相等，即 \(-\dfrac{a}{2}=3\)，解得 \(a=-6\)，故选 B.
\end{solution}

\begin{problem}
函数 \(y = \tan \left(\frac{1}{2}x - \frac{\pi}{3}\right)\) 在一个周期内的图象是（\(\quad\)）

\choices
    {\choicebitmap[width=0.15\paperwidth]{img/1997/national_science/q3_optA_clean.png}}
    {\choicebitmap[width=0.15\paperwidth]{img/1997/national_science/q3_optB_clean.png}}
    {\choicebitmap[width=0.15\paperwidth]{img/1997/national_science/q3_optC_clean.png}}
    {\choicebitmap[width=0.15\paperwidth]{img/1997/national_science/q3_optD_clean.png}}
\end{problem}

\begin{answer}
A
\end{answer}

\begin{solution}
令 \(u=\dfrac{x}{2}-\dfrac{\pi}{3}\). 正切函数的周期为 \(\pi\)，所以原函数的最小正周期为 \(\dfrac{\pi}{1/2}=2\pi\). 在相邻两条渐近线之间，\(u\in\left(-\dfrac{\pi}{2},\dfrac{\pi}{2}\right)\)，对应 \(x\in\left(-\dfrac{\pi}{3},\dfrac{5\pi}{3}\right)\)，函数单调递增；又 \(u=0\) 时 \(x=\dfrac{2\pi}{3}\)，图象经过 \(\left(\dfrac{2\pi}{3},0\right)\). 这与选项 A 的图象相符，故选 A.
\end{solution}

\begin{problem}
已知三棱锥 \(D - ABC\) 的三个侧面与底面全等，且 \(AB = AC = \sqrt{3}\)，\(BC = 2\)，则以 \(BC\) 为棱，以面 \(BCD\) 与面 \(BCA\) 为面的二面角的大小是（\(\quad\)）

\choices
    {\(\arccos \frac{\sqrt{3}}{3}\)}
    {\(\arccos \frac{1}{3}\)}
    {\(\frac{\pi}{2}\)}
    {\(\frac{2\pi}{3}\)}
\end{problem}

\begin{answer}
C
\end{answer}

\begin{solution}
因为三个侧面都与底面全等，面 \(BCD\) 的三边仍为 \(2\)，\(\sqrt{3}\)，\(\sqrt{3}\)，故 \(BD=CD=\sqrt{3}\). 面 \(ABD\) 与底面全等，且 \(AB=BD=\sqrt{3}\)，所以 \(AD=2\). 取坐标 \(B=(-1,0,0)\)，\(C=(1,0,0)\)，\(A=(0,\sqrt{2},0)\). 设 \(D=(u,v,w)\). 由

\[
DB^{2}-DC^{2}=0
\]

得 \(u=0\)，再由 \(DB=\sqrt{3}\) 得 \(v^{2}+w^{2}=2\). 由 \(DA=2\) 得

\[
(v-\sqrt{2})^{2}+w^{2}=4
\]
将该式与 \(v^{2}+w^{2}=2\) 相减，得 \(2-2\sqrt{2}v=2\)，即 \(v=0\)，从而 \(w^{2}=2\). 取 \(w=\sqrt{2}\)；取负值只对应关于底面 \(ABC\) 的对称位置，不影响二面角. 此时 \(D=(0,0,\sqrt{2})\)，平面 \(BCA\) 的方程为 \(z=0\)，平面 \(BCD\) 的方程为 \(y=0\). 两平面的法向量分别为 \((0,0,1)\) 与 \((0,1,0)\)，相互垂直，因此所求二面角为 \(\dfrac{\pi}{2}\)，故选 C.
\end{solution}

\begin{problem}
函数 \(y = \sin \left(\frac{\pi}{3} - 2x\right) + \cos 2x\) 的最小正周期是（\(\quad\)）

\choices
    {\(\frac{\pi}{2}\)}
    {\(\pi\)}
    {\(2\pi\)}
    {\(4\pi\)}
\end{problem}

\begin{answer}
B
\end{answer}

\begin{solution}
利用正弦差角公式，

\[
\sin\left(\frac{\pi}{3}-2x\right)+\cos 2x=\left(\frac{\sqrt{3}}{2}+1\right)\cos 2x-\frac{1}{2}\sin 2x
\]

右端是非零的 \(\cos 2x\) 与 \(\sin 2x\) 的线性组合，角的频率为 \(2\)，故最小正周期为 \(\dfrac{2\pi}{2}=\pi\)，选 B.
\end{solution}

\begin{problem}
满足 \(\arccos(1 - x)\geqslant \arccos x\) 的 \(x\) 的取值范围是（\(\quad\)）

\choices
    {\(\left[-1, -\frac{1}{2}\right]\)}
    {\(\left[-\frac{1}{2},0\right]\)}
    {\(\left[0,\frac{1}{2}\right]\)}
    {\(\left[\frac{1}{2},1\right]\)}
\end{problem}

\begin{answer}
D
\end{answer}

\begin{solution}
由两个反余弦式有意义，必须同时满足 \(-1\leqslant 1-x\leqslant1\) 和 \(-1\leqslant x\leqslant1\)，故 \(x\in[0,1]\). 反余弦函数在 \([-1,1]\) 上单调递减，所以

\[
\arccos(1-x)\geqslant\arccos x\Longleftrightarrow 1-x\leqslant x\Longleftrightarrow x\geqslant\frac12
\]

因此 \(x\in\left[\dfrac12,1\right]\)，故选 D.
\end{solution}

\begin{problem}
将 \(y = 2^x\) 的图象\fillinblank{}，再作关于直线 \(y = x\) 对称的图象，可得到函数 \(y = \log_{2}(x + 1)\) 的图象（\(\quad\)）

\choices
    {先向左平行移动 \(1\) 个单位}
    {先向右平行移动 \(1\) 个单位}
    {先向上平行移动 \(1\) 个单位}
    {先向下平行移动 \(1\) 个单位}
\end{problem}

\begin{answer}
D
\end{answer}

\begin{solution}
先把 \(y=2^{x}\) 的图象向下平移 \(1\) 个单位，得到 \(y=2^{x}-1\). 关于 \(y=x\) 对称后，交换 \(x\) 与 \(y\)，有 \(x=2^{y}-1\)，即 \(y=\log_{2}(x+1)\). 因此应先向下平移 \(1\) 个单位，选 D.
\end{solution}

\begin{problem}
长方体一个顶点上三条棱的长分别是 \(3\)，\(4\)，\(5\)，且它的八个顶点都在同一个球面上，这个球的表面积是（\(\quad\)）

\choices
    {\(20\sqrt{2}\pi\)}
    {\(25\sqrt{2}\pi\)}
    {\(50\pi\)}
    {\(200\pi\)}
\end{problem}

\begin{answer}
C
\end{answer}

\begin{solution}
长方体的体对角线长为 \(d=\sqrt{3^{2}+4^{2}+5^{2}}=5\sqrt{2}\). 球心是长方体体对角线的中点，所以球半径 \(R=\dfrac{d}{2}=\dfrac{5\sqrt{2}}{2}\). 因而球的表面积为 \(4\pi R^{2}=4\pi\cdot\dfrac{50}{4}=50\pi\)，故选 C.
\end{solution}

\begin{problem}
曲线的参数方程是 \(\begin{cases}
x = 1 - \frac{1}{t},\\
y = 1 - t^2,
\end{cases}\)（\(t\) 是参数，\(t \neq 0\)），它的普通方程是（\(\quad\)）

\choices
    {\((x - 1)^{2}(y - 1) = 1\)}
    {\(y = \frac{x(x - 2)}{(1 - x)^2}\)}
    {\(y = \frac{1}{(1 - x)^2} - 1\)}
    {\(y = \frac{x}{1 - x^2} + 1\)}
\end{problem}

\begin{answer}
B
\end{answer}

\begin{solution}
由 \(x=1-\dfrac1t\) 且 \(t\ne0\)，得 \(1-x=\dfrac1t\)，所以 \(t=\dfrac1{1-x}\)，并且 \(x\ne1\). 代入 \(y=1-t^{2}\)，得

\[
y=1-\frac{1}{(1-x)^{2}}=\frac{(1-x)^{2}-1}{(1-x)^{2}}=\frac{x(x-2)}{(1-x)^{2}}
\]

故普通方程为选项 B.
\end{solution}

\begin{problem}
函数 \(y = \cos^2 x - 3\cos x + 2\) 的最小值为（\(\quad\)）

\choices
    {\(2\)}
    {\(0\)}
    {\(-\frac{1}{4}\)}
    {\(6\)}
\end{problem}

\begin{answer}
B
\end{answer}

\begin{solution}
令 \(u=\cos x\)，则 \(-1\leqslant u\leqslant1\)，且

\[
y=u^{2}-3u+2=(u-1)(u-2)
\]

在 \([-1,1]\) 上，\(u-1\leqslant0\)，\(u-2<0\)，且 \(u^{2}-3u+2\) 的导数 \(2u-3<0\)，所以函数随 \(u\) 增大而减小，最小值在 \(u=1\) 时取得，等于 \(0\). 故选 B.
\end{solution}

\begin{problem}
椭圆 \(C\) 与椭圆 \(\frac{(x - 3)^{2}}{9} + \frac{(y - 2)^{2}}{4} = 1\) 关于直线 \(x + y = 0\) 对称，椭圆 \(C\) 的方程是（\(\quad\)）

\choices
    {\(\frac{(x + 2)^2}{4} + \frac{(y + 3)^2}{9} = 1\)}
    {\(\frac{(x - 2)^2}{9} + \frac{(y - 3)^2}{4} = 1\)}
    {\(\frac{(x + 2)^2}{9} + \frac{(y + 3)^2}{4} = 1\)}
    {\(\frac{(x - 2)^2}{4} + \frac{(y - 3)^2}{9} = 1\)}
\end{problem}

\begin{answer}
A
\end{answer}

\begin{solution}
关于直线 \(x+y=0\)，即 \(y=-x\) 对称的点变换为 \((x,y)\mapsto(-y,-x)\). 设椭圆 \(C\) 上的点为 \((X,Y)\)，其对称点 \((-Y,-X)\) 在已知椭圆上，因此

\[
\frac{(-Y-3)^{2}}{9}+\frac{(-X-2)^{2}}{4}=1
\]

化简得

\[
\frac{(X+2)^{2}}{4}+\frac{(Y+3)^{2}}{9}=1
\]

故选 A.
\end{solution}

\begin{problem}
圆台上，下底面积分别为 \(\pi\)，\(4\pi\)，侧面积为 \(6\pi\)，这个圆台体积是（\(\quad\)）

\choices
    {\(\frac{2\sqrt{3}\pi}{3}\)}
    {\(2\sqrt{3}\pi\)}
    {\(\frac{7\sqrt{3}\pi}{6}\)}
    {\(\frac{7\sqrt{3}\pi}{3}\)}
\end{problem}

\begin{answer}
D
\end{answer}

\begin{solution}
设圆台上下底半径分别为 \(r\)，\(R\)，母线长为 \(l\)，高为 \(h\). 由底面积得 \(r=1\)，\(R=2\). 侧面积满足 \(\pi(R+r)l=6\pi\)，所以 \(l=2\). 于是

\[
h=\sqrt{l^{2}-(R-r)^{2}}=\sqrt{4-1}=\sqrt{3}
\]

圆台体积为

\[
V=\frac{\pi h}{3}(R^{2}+Rr+r^{2})=\frac{\pi\sqrt{3}}{3}(4+2+1)=\frac{7\sqrt{3}\pi}{3}
\]

故选 D.
\end{solution}

\begin{problem}
定义在区间 \((-\infty, +\infty)\) 的奇函数 \(f(x)\) 为增函数，偶函数 \(g(x)\) 在区间 \([0, +\infty)\) 的图象与 \(f(x)\) 的图象重合，设 \(a > b > 0\)，给出下列不等式：

\begin{circlelist}
\item \(f(b)-f(-a)>g(a)-g(-b)\)；
\item \(f(b)-f(-a)<g(a)-g(-b)\)；
\item \(f(a)-f(-b)>g(b)-g(-a)\)；
\item \(f(a)-f(-b)<g(b)-g(-a)\).
\end{circlelist}

其中成立的是（\(\quad\)）

\choices
    {\(\circled{1}\)与\(\circled{4}\)}
    {\(\circled{2}\)与\(\circled{3}\)}
    {\(\circled{1}\)与\(\circled{3}\)}
    {\(\circled{2}\)与\(\circled{4}\)}
\end{problem}

\begin{answer}
C
\end{answer}

\begin{solution}
由 \(f\) 为奇函数、\(g\) 为偶函数，且 \(g\) 在非负半轴上与 \(f\) 重合，得 \(f(-a)=-f(a)\)，\(f(-b)=-f(b)\)，\(g(a)=f(a)\)，\(g(-a)=f(a)\)，\(g(b)=f(b)\)，\(g(-b)=f(b)\). 又因为 \(a>b>0\) 且 \(f\) 递增，所以 \(f(a)>f(b)>f(0)=0\). 因此第一个不等式两边分别为 \(f(a)+f(b)\) 与 \(f(a)-f(b)\)，第一个不等式成立，第二个不等式不成立；第三个不等式两边分别为 \(f(a)+f(b)\) 与 \(f(b)-f(a)\)，第三个不等式成立，第四个不等式不成立. 故成立的是 \(\circled{1}\) 与 \(\circled{3}\)，选 C.
\end{solution}

\begin{problem}
不等式组 \(\begin{cases}
x>0,\\
\frac{3-x}{3+x} > \left|\frac{2-x}{2+x}\right|
\end{cases}\) 的解集是（\(\quad\)）

\choices
    {\(\{x \mid 0 < x < 2\}\)}
    {\(\{x \mid 0 < x < 2.5\}\)}
    {\(\{x \mid 0 < x < \sqrt{6}\}\)}
    {\(\{x \mid 0 < x < 3\}\)}
\end{problem}

\begin{answer}
C
\end{answer}

\begin{solution}
由 \(x>0\) 可知 \(2+x>0\)，\(3+x>0\). 原不等式右端非负，所以左端必须为正，从而 \(0<x<3\). 在此范围内两边均非负，可以平方，得到

\[
\left(\frac{3-x}{3+x}\right)^{2}>\left(\frac{2-x}{2+x}\right)^{2}
\]

由于分母为正，令 \(A=(3-x)(2+x)\)，\(B=(2-x)(3+x)\)，则上式等价于 \(A^{2}>B^{2}\). 计算得 \(A-B=2x>0\)，\(A+B=12-2x^{2}\)，所以 \(A^{2}-B^{2}>0\) 等价于 \(x^{2}<6\). 与 \(0<x<3\) 合并，得 \(0<x<\sqrt{6}\)，故选 C.
\end{solution}

\begin{problem}
四面体的顶点和各棱中点共 \(10\) 个点，在其中取 \(4\) 个不共面的点，不同的取法共有（\(\quad\)）

\choices
    {\(150\) 种}
    {\(147\) 种}
    {\(144\) 种}
    {\(141\) 种}
\end{problem}

\begin{answer}
D
\end{answer}

\begin{solution}
从 \(10\) 个点中任取 \(4\) 个共有 \(\mathrm C_{10}^{4}=210\) 种. 先数四点共面的取法. 每个面的三个顶点和三条棱的中点共 \(6\) 个点，从一个面内取 \(4\) 个有 \(\mathrm C_{6}^{4}=15\) 种；四个面两两的交集只有一条棱上的 \(3\) 个已知点，不可能有四点同时被重复计数，所以面内共面取法共有 \(4\mathrm C_{6}^{4}=60\) 种.

设四面体顶点的仿射坐标为 \(A=(1,0,0)\)，\(B=(0,1,0)\)，\(C=(0,0,1)\)，\(D=(0,0,0)\)，并记棱 \(UV\) 的中点为 \(M_{UV}\). 四点 \(P_i=(x_i,y_i,z_i)\) 共面的充要条件是

\[
\det\begin{pmatrix}
x_2-x_1&y_2-y_1&z_2-z_1\\
x_3-x_1&y_3-y_1&z_3-z_1\\
x_4-x_1&y_4-y_1&z_4-z_1
\end{pmatrix}=0
\]

每个面内有 \(3\) 个顶点和 \(3\) 个棱中点. 取 \(3\) 个顶点时，这三点确定一个面，第四点若要共面，只能是该面三条棱中的一个中点；因此这类共面组已经计入上述 \(60\) 组. 取 \(1\) 个顶点时，以顶点 \(A\) 为例，六个棱中点记为 \(M_{AB}\)，\(M_{AC}\)，\(M_{AD}\)，\(M_{BC}\)，\(M_{BD}\)，\(M_{CD}\). 逐一代入上面的行列式，行列式为 \(0\) 的三组中点恰为

\(\{M_{AB},M_{AC},M_{BC}\}\)，\(\{M_{AB},M_{AD},M_{BD}\}\)，\(\{M_{AC},M_{AD},M_{CD}\}\).

例如第一组与 \(A\) 同在面 \(ABC\) 内，第二组与 \(A\) 同在面 \(ABD\) 内，第三组与 \(A\) 同在面 \(ACD\) 内；其余 \(17\) 组的行列式均不为 \(0\). 其他顶点的情形由四面体顶点置换对称. 所以含 \(3\) 个或 \(1\) 个顶点的共面组均已计入 \(60\) 组，面外情形只需考察含 \(2\) 个或 \(0\) 个顶点的情形.

先考察含 \(2\) 个顶点的情形. 以顶点 \(A\)，\(B\) 为例，六个棱中点中任取两个，共有 \(\mathrm C_6^2=15\) 种. 对向量 \(\overrightarrow{AB}=(-1,1,0)\) 和任意两个中点 \(M\)，\(N\)，共面条件是

\[
\det\begin{pmatrix}
-1&x_M-1&x_N-1\\
1&y_M&y_N\\
0&z_M&z_N
\end{pmatrix}=0
\]

六个面内的取法为 \(\{M_{AB},M_{AC}\}\)、\(\{M_{AB},M_{BC}\}\)、\(\{M_{AC},M_{BC}\}\) 以及 \(\{M_{AB},M_{AD}\}\)、\(\{M_{AB},M_{BD}\}\)、\(\{M_{AD},M_{BD}\}\)，对应面 \(ABC\) 或面 \(ABD\)，正好是已计入面内的 \(6\) 种. 其余 \(9\) 种逐一代入上式：

\[
\begin{gathered}
\det=0\Longleftrightarrow\{M,N\}=\{M_{AB},M_{CD}\},\\
\det=-\dfrac14\quad\text{对应}\quad\{M_{AC},M_{AD}\},\{M_{AC},M_{BD}\},\{M_{AC},M_{CD}\},\{M_{BC},M_{BD}\},\{M_{BC},M_{CD}\},\\
\det=\dfrac14\quad\text{对应}\quad\{M_{AD},M_{BC}\},\{M_{AD},M_{CD}\},\{M_{BD},M_{CD}\}.
\end{gathered}
\]

因此，除面内取法外，每条棱的两个端点与该棱及其对棱的中点还给出 \(1\) 组额外共面点集. 六条棱共得 \(6\) 组：

\(\{A,B,M_{AB},M_{CD}\}\)，\(\ \{C,D,M_{AB},M_{CD}\}\)，\(\ \{A,C,M_{AC},M_{BD}\}\)，\(\ \{B,D,M_{AC},M_{BD}\}\)

\(\{A,D,M_{AD},M_{BC}\}\)，\(\ \{B,C,M_{AD},M_{BC}\}\).

再考察 \(0\) 个顶点的情形. 四个中点的集合等价于从六个中点中删去两个. 将剩余四个中点代入上面的行列式判据，行列式为 \(0\) 的情形恰好是删去一对对棱，即

\(\{M_{AB},M_{AC},M_{BD},M_{CD}\}\)，\(\{M_{AB},M_{AD},M_{BC},M_{CD}\}\)，\(\{M_{AC},M_{AD},M_{BC},M_{BD}\}\).

其余 \(12\) 种删法代入后行列式均为 \(\pm\dfrac18\)，不为 \(0\). 上述 \(3\) 组均不在单一面内，且分类已覆盖顶点数为 \(0\)，\(1\)，\(2\)，\(3\)，\(4\) 的全部可能. 因而共面取法有 \(60+6+3=69\) 种，不共面的取法为 \(210-69=141\) 种，故选 D.
\end{solution}

\section{填空题}

\begin{problem}
已知 \(\left(\frac{a}{x} - \sqrt{\frac{x}{2}}\right)^{9}\) 的展开式中 \(x^{3}\) 的系数为 \(\frac{9}{4}\)，常数 \(a\) 的值为\fillinblank{}.
\end{problem}

\begin{answer}
\(4\)
\end{answer}

\begin{solution}
展开式通项中取 \(k\) 个 \(\dfrac{a}{x}\)，取 \(9-k\) 个 \(-\sqrt{\dfrac{x}{2}}\)，得

\[
\mathrm C_{9}^{k}\left(\frac{a}{x}\right)^{k}\left(-\sqrt{\frac{x}{2}}\right)^{9-k}
\]

其中 \(x\) 的指数为 \(-k+\dfrac{9-k}{2}=\dfrac{9-3k}{2}\). 要得到 \(x^{3}\)，必须有 \(\dfrac{9-3k}{2}=3\)，解得 \(k=1\). 此时 \(x^{3}\) 的系数为 \(\mathrm C_{9}^{1}a\left(\dfrac{1}{\sqrt{2}}\right)^{8}=\dfrac{9a}{16}\). 由 \(\dfrac{9a}{16}=\dfrac94\)，得 \(a=4\).
\end{solution}

\begin{problem}
已知直线的极坐标方程为 \(\rho \sin \left(\theta + \frac{\pi}{4}\right) = \frac{\sqrt{2}}{2}\)，则极点到该直线的距离是\fillinblank{}.
\end{problem}

\begin{answer}
\(\frac{\sqrt{2}}{2}\)
\end{answer}

\begin{solution}
由 \(x=\rho\cos\theta\)，\(y=\rho\sin\theta\)，得

\[
\rho\sin\left(\theta+\frac{\pi}{4}\right)=\frac{x+y}{\sqrt{2}}
\]

所以极坐标方程化为 \(x+y=1\)，即 \(x+y-1=0\). 极点就是原点，其到该直线的距离为 \(\dfrac{|0+0-1|}{\sqrt{1^{2}+1^{2}}}=\dfrac{\sqrt{2}}{2}\).
\end{solution}

\begin{problem}
\(\frac{\sin 7^{\circ} + \cos 15^{\circ} \sin 8^{\circ}}{\cos 7^{\circ} - \sin 15^{\circ} \sin 8^{\circ}}\) 的值为\fillinblank{}.
\end{problem}

\begin{answer}
\(2 - \sqrt{3}\)
\end{answer}

\begin{solution}
由积化和差公式，

\[
\cos15^{\circ}\sin8^{\circ}=\frac12\left(\sin23^{\circ}-\sin7^{\circ}\right)
\]

所以分子为 \(\dfrac12(\sin7^{\circ}+\sin23^{\circ})=\sin15^{\circ}\cos8^{\circ}\). 同理，\(\sin15^{\circ}\sin8^{\circ}=\dfrac12(\cos7^{\circ}-\cos23^{\circ})\)，故分母为 \(\dfrac12(\cos7^{\circ}+\cos23^{\circ})=\cos15^{\circ}\cos8^{\circ}\). 因而原式等于 \(\tan15^{\circ}=2-\sqrt{3}\).
\end{solution}

\begin{problem}
已知 \(m\)，\(l\) 是直线，\(\alpha\)，\(\beta\) 是平面，给出下列命题：

\begin{circlelist}
\item 若 \(l\) 垂直于 \(\alpha\) 内的两条相交直线，则 \(l \perp \alpha\)；
\item 若 \(l\) 平行于 \(\alpha\)，则 \(l\) 平行于 \(\alpha\) 内的所有直线；
\item 若 \(m \subset \alpha\)，\(l \subset \beta\)，且 \(l \perp m\)，则 \(\alpha \perp \beta\)；
\item 若 \(l \subset \beta\)，且 \(l \perp \alpha\)，则 \(\alpha \perp \beta\)；
\item 若 \(m \subset \alpha\)，\(l \subset \beta\)，且 \(\alpha \parallel \beta\)，则 \(m \parallel l\).
\end{circlelist}

其中正确的命题的序号是\fillinblank{}.
\end{problem}

\begin{answer}
\(\circled{1}\)，\(\circled{4}\)
\end{answer}

\begin{solution}
命题 \(\circled{1}\) 是平面垂直的判定定理：若一条直线垂直于一个平面内的两条相交直线，则这条直线垂直于该平面，所以 \(\circled{1}\) 正确.

命题 \(\circled{2}\) 不正确. 例如取 \(\alpha\) 为平面 \(z=0\)，取直线 \(l=\{(t,0,1)\mid t\in\R\}\). 则 \(l\parallel\alpha\)，但平面 \(\alpha\) 内的直线 \(n=\{(0,t,0)\mid t\in\R\}\) 的方向与 \(l\) 不同，故 \(l\) 不平行于 \(n\)，\(\circled{2}\) 错误.

命题 \(\circled{3}\) 不正确. 取 \(\alpha\) 为平面 \(z=0\)，\(\beta\) 为平面 \(x+z=0\)，取 \(m=\{(t,0,0)\mid t\in\R\}\)，\(l=\{(0,t,0)\mid t\in\R\}\). 两直线在原点相交且 \(l\perp m\)，但 \(\alpha\) 与 \(\beta\) 的法向量分别为 \((0,0,1)\) 和 \((1,0,1)\)，不垂直，所以不能推出 \(\alpha\perp\beta\)，\(\circled{3}\) 错误.

命题 \(\circled{4}\) 也是平面垂直的判定定理：若一个平面经过另一个平面的垂线，则这两个平面垂直，所以 \(\circled{4}\) 正确.

命题 \(\circled{5}\) 不正确. 取 \(\alpha\) 为平面 \(z=0\)，\(\beta\) 为平面 \(z=1\)，取 \(m\) 为 \(\alpha\) 内平行于 \(x\) 轴的直线，取 \(l\) 为 \(\beta\) 内平行于 \(y\) 轴的直线. 此时 \(\alpha\parallel\beta\)，但 \(m\) 与 \(l\) 的方向不同，故 \(m\) 与 \(l\) 不平行，\(\circled{5}\) 错误.

因此正确的命题序号为 \(\circled{1}\)，\(\circled{4}\).
\end{solution}

\section{解答题}

\begin{problem}
已知复数 \(z = \frac{\sqrt{3}}{2} - \frac{1}{2}\i\)，\(\omega = \frac{\sqrt{2}}{2} + \frac{\sqrt{2}}{2}\i\). 复数 \(\overline{z\omega}\)，\(z^{2}\omega^{3}\) 在复数平面上所对应的点分别为 \(P\)，\(Q\). 证明 \(\triangle OPQ\) 是等腰直角三角形（其中 \(O\) 为原点）.
\end{problem}

\begin{solution}
由 \(z=\cos\left(-\frac{\pi}{6}\right)+\i\sin\left(-\frac{\pi}{6}\right)\)，\(\omega=\cos\frac{\pi}{4}+\i\sin\frac{\pi}{4}\)，得 \(|z|=|\omega|=1\)，\(\arg z=-\frac{\pi}{6}\)，\(\arg\omega=\frac{\pi}{4}\). 因而

\(\overline{z\omega}=\cos\left(-\frac{\pi}{12}\right)+\i\sin\left(-\frac{\pi}{12}\right)\)，\(z^{2}\omega^{3}=\cos\frac{5\pi}{12}+\i\sin\frac{5\pi}{12}\).

所以 \(|\overline{z\omega}|=|z^{2}\omega^{3}|=1\)，从而 \(OP=OQ=1\). 两个辐角之差为

\[
\frac{5\pi}{12}-\left(-\frac{\pi}{12}\right)=\frac{\pi}{2}
\]

所以 \(\angle POQ=\frac{\pi}{2}\). 因此 \(OP=OQ\) 且 \(\angle POQ=90^\circ\)，\(\triangle OPQ\) 是等腰直角三角形.
\end{solution}

\begin{problem}
已知数列 \(\{a_{n}\}\)，\(\{b_{n}\}\) 都是由正数组成的等比数列，公比分别为 \(p\)，\(q\)，其中 \(p > q\)，且 \(p \neq 1\)，\(q \neq 1\). 设 \(c_{n} = a_{n} + b_{n}\)，\(S_{n}\) 为数列 \(\{c_{n}\}\) 的前 \(n\) 项和. 求 \(\displaystyle\lim_{n \to \infty} \frac{S_{n}}{S_{n-1}}\).
\end{problem}

\begin{solution}
因为两个数列都由正数组成，所以 \(a_{1}>0\)，\(b_{1}>0\)，且 \(p>q>0\). 由等比数列求和公式，得
\[
S_{n}=\frac{a_{1}(p^{n}-1)}{p-1}+\frac{b_{1}(q^{n}-1)}{q-1}
\]
先设 \(p>1\). 因为 \(0<q/p<1\)，所以
\[
\frac{S_{n}}{p^{n}}=\frac{a_{1}(1-p^{-n})}{p-1}+\frac{b_{1}}{q-1}\left[\left(\frac{q}{p}\right)^{n}-p^{-n}\right]\longrightarrow\frac{a_{1}}{p-1}>0
\]
把 \(n\) 换成 \(n-1\)，同样有
\[
\frac{S_{n-1}}{p^{n-1}}\longrightarrow\frac{a_{1}}{p-1}>0
\]
因此
\[
\frac{S_{n}}{S_{n-1}}=p\frac{S_{n}/p^{n}}{S_{n-1}/p^{n-1}}\longrightarrow p
\]
再设 \(p<1\). 由 \(0<q<p<1\)，得
\[
S_{n}\longrightarrow\frac{a_{1}}{1-p}+\frac{b_{1}}{1-q}>0
\]
并且 \(S_{n-1}\) 也趋于同一个正数. 于是分母不趋于 \(0\)，且
\[
\frac{S_{n}}{S_{n-1}}\longrightarrow1
\]
因此
\[
\lim_{n\to\infty}\frac{S_{n}}{S_{n-1}}=
\begin{cases}
p,&p>1\\
1,&0<p<1
\end{cases}
\]
\end{solution}

\begin{problem}
甲，乙两地相距 \(S\text{ 千米}\)，汽车从甲地匀速行驶到乙地，速度不得超过 \(c\text{ 千米/时}\). 已知汽车每小时的运输成本（以元为单位）由可变部分和固定部分组成：可变部分与速度 \(v\text{ 千米/时}\) 的平方成正比，比例系数为 \(b\)；固定部分为 \(a\text{ 元}\).

\begin{enumerate}
\item 把全程运输成本 \(y\text{ 元}\) 表示为速度 \(v\text{ 千米/时}\) 的函数，并指出这个函数的定义域；

\item 为了使全程运输成本最小，汽车应以多大速度行驶？
\end{enumerate}
\end{problem}

\begin{solution}
\begin{enumerate}
\item 由题意 \(a>0\)，\(b>0\). 汽车以速度 \(v\) 行驶全程所需时间为 \(S/v\)，每小时成本为 \(a+bv^{2}\)，
所以 \(y=\frac{S}{v}(a+bv^{2})=S\left(\frac{a}{v}+bv\right)\)，函数定义域为 \(0<v\leqslant c\).
\item 在定义域内，
\(y'(v)=S\left(b-\frac{a}{v^{2}}\right)\)，\(y''(v)=\frac{2aS}{v^{3}}>0\).
令 \(y'(v)=0\)，得临界速度 \(v_{0}=\sqrt{a/b}\). 由于 \(y''(v)>0\)，函数在临界点处取得极小值. 若 \(v_{0}\leqslant c\)，最优速度为 \(\sqrt{a/b}\)；若 \(v_{0}>c\)，则 \(y'(v)<0\) 在 \((0,c]\) 上成立，函数递减，最优速度为 \(c\).
\end{enumerate}
\end{solution}

\begin{problem}
如图，在正方体 \(ABCD - A_{1}B_{1}C_{1}D_{1}\) 中，\(E\)，\(F\) 分别是 \(BB_{1}\)，\(CD\) 的中点.

\begin{enumerate}
\item 证明：\(AD \perp D_{1}F\)；

\item 求 \(AE\) 与 \(D_{1}F\) 所成的角；

\item 证明：面 \(AED\) \(\perp\) 面 \(A_{1}FD_{1}\)；

\item 设 \(AA_{1} = 2\)，求三棱锥 \(F - A_{1}ED_{1}\) 的体积.
\end{enumerate}

\bitmapfigure[width=0.33\linewidth]{img/1997/national_science/q23_fig1.png}
\end{problem}

\begin{solution}
设正方体棱长为 \(a\)，建立直角坐标系，使 \(A=(0,0,0)\)，\(B=(a,0,0)\)，\(C=(a,a,0)\)，\(D=(0,a,0)\)，且 \(A_{1}=(0,0,a)\)，\(B_{1}=(a,0,a)\)，\(C_{1}=(a,a,a)\)，\(D_{1}=(0,a,a)\).
于是 \(E=(a,0,a/2)\)，\(F=(a/2,a,0)\)，从而 \(\overrightarrow{AD}=(0,a,0)\)，\(\overrightarrow{D_{1}F}=(a/2,0,-a)\)，\(\overrightarrow{AE}=(a,0,a/2)\).
\begin{enumerate}
\item \(\overrightarrow{AD}\cdot\overrightarrow{D_{1}F}=0\)，故 \(AD\perp D_{1}F\).
\item \(\overrightarrow{AE}\cdot\overrightarrow{D_{1}F}=a\cdot\dfrac{a}{2}+\dfrac{a}{2}\cdot(-a)=0\)，故 \(AE\) 与 \(D_{1}F\) 所成的角为 \(\dfrac{\pi}{2}\).
\item \(AD\) 与 \(AE\) 是平面 \(AED\) 内的两条相交直线，且由上面的数量积计算可知 \(D_{1}F\perp AD\)，\(D_{1}F\perp AE\). 因而 \(\overrightarrow{D_{1}F}\) 是平面 \(AED\) 的法向量. 又平面 \(AED\) 的方程为 \(z=x/2\)，直线 \(D_{1}F\) 上的点可写为 \((ta/2,a,a(1-t))\)，取 \(t=4/5\) 时满足该方程，所以 \(D_{1}F\) 与平面 \(AED\) 相交，从而 \(D_{1}F\perp\) 平面 \(AED\). 由于 \(D_{1}F\) 在平面 \(A_{1}FD_{1}\) 内，故面 \(AED\perp\) 面 \(A_{1}FD_{1}\).
\item 取 \(F\) 为顶点，三棱锥的三个棱向量为 \(\overrightarrow{FA_{1}}=\left(-\frac{a}{2},-a,a\right)\)，\(\overrightarrow{FE}=\left(\frac{a}{2},-a,\frac{a}{2}\right)\)，\(\overrightarrow{FD_{1}}=\left(-\frac{a}{2},0,a\right)\).
因此
\[
V=\frac{1}{6}\left|\det\begin{pmatrix}
-\frac{a}{2}&\frac{a}{2}&-\frac{a}{2}\\
-a&-a&0\\
a&\frac{a}{2}&a
\end{pmatrix}\right|
\]
展开行列式得
\[
\det\begin{pmatrix}
-\frac{a}{2}&\frac{a}{2}&-\frac{a}{2}\\
-a&-a&0\\
a&\frac{a}{2}&a
\end{pmatrix}
=a^{3}\left[\left(-\frac{1}{2}\right)(-1)-\frac{1}{2}(-1)+\left(-\frac{1}{2}\right)\frac{1}{2}\right]
=\frac{3a^{3}}{4}
\]
所以 \(V=\frac{1}{6}\cdot\frac{3a^{3}}{4}=\frac{a^{3}}{8}\).
由 \(AA_{1}=a=2\)，得 \(V=1\).
\end{enumerate}
\end{solution}

\begin{problem}
设二次函数 \(f(x) = ax^{2} + bx + c\)（\(a > 0\)），方程 \(f(x) - x = 0\) 的两个根 \(x_{1}\)，\(x_{2}\) 满足 \(0 < x_{1} < x_{2} < \frac{1}{a}\).

\begin{enumerate}
\item 当 \(x \in (0, x_{1})\) 时，证明 \(x < f(x) < x_{1}\)；

\item 设函数 \(f(x)\) 的图象关于直线 \(x = x_{0}\) 对称，证明 \(x_{0} < \frac{x_{1}}{2}\).
\end{enumerate}
\end{problem}

\begin{solution}
由根与因式的关系，\(f(x)-x=a(x-x_{1})(x-x_{2})\).
\begin{enumerate}
\item 当 \(0<x<x_{1}\) 时，\(x-x_{1}<0\)，\(x-x_{2}<0\)，
所以 \(f(x)-x=a(x-x_{1})(x-x_{2})>0\). 又 \(f(x)-x_{1}=(x-x_{1})\bigl(1+a(x-x_{2})\bigr)\).
因为 \(x<x_{1}\) 且 \(x_{2}<1/a\)，所以 \(x_{2}-x<x_{2}<1/a\)，从而 \(1+a(x-x_{2})=1-a(x_{2}-x)>0\). 因此 \(f(x)-x_{1}<0\)，即 \(x<f(x)<x_{1}\).
\item 比较
\(a(x-x_{1})(x-x_{2})=ax^{2}+(b-1)x+c\)
的 \(x\) 项系数，得 \(b-1=-a(x_{1}+x_{2})\)，即 \(b=1-a(x_{1}+x_{2})\). 抛物线的对称轴为 \(x=x_{0}\)，故
\(x_{0}=-\frac{b}{2a}=\frac{x_{1}+x_{2}}{2}-\frac{1}{2a}\).
由 \(x_{2}<1/a\)，得
\(x_{0}<\frac{x_{1}+1/a}{2}-\frac{1}{2a}=\frac{x_{1}}{2}\).
\end{enumerate}
\end{solution}

\begin{problem}
设圆满足：

\begin{circlelist}
\item 截 \(y\) 轴所得弦长为 \(2\)；
\item 被 \(x\) 轴分成两段圆弧，其弧长的比为 \(3:1\).
\end{circlelist}

在满足上述两条件的所有圆中，求圆心到直线 \(l\colon x - 2y = 0\) 的距离最小的圆的方程.
\end{problem}

\begin{solution}
设圆心为 \((h,k)\)，半径为 \(r\). 设圆与 \(y\) 轴的弦中点为 \(H\). 由圆心到弦的距离为 \(|h|\)，弦长为 \(2\)，得
\(2\sqrt{r^{2}-h^{2}}=2\)，\(r^{2}-h^{2}=1\).
圆被 \(x\) 轴截成两段圆弧，弧长比等于所对圆心角之比，所以两圆心角为 \(3\pi/2\) 与 \(\pi/2\). 取劣弧对应的圆心角 \(\pi/2\)，其弦到圆心的距离为
\[
|k|=r\cos\frac{\pi}{4}=\frac{r}{\sqrt2}
\]
故 \(r^{2}=2k^{2}\)，再结合 \(r^{2}-h^{2}=1\)，得到
\(h^{2}=2k^{2}-1\)，\(|k|\geqslant\frac1{\sqrt2}\).
圆心到 \(l:x-2y=0\) 的距离为 \(d=\frac{|h-2k|}{\sqrt5}\). 令 \(u=|k|\)，则 \(|h|=\sqrt{2u^{2}-1}\). 对固定的 \(u\)，若 \(h\)，\(k\) 同号，则
\[
|h-2k|=2u-\sqrt{2u^{2}-1};
\]
若 \(h\)，\(k\) 异号，则
\[
|h-2k|=2u+\sqrt{2u^{2}-1}>2u-\sqrt{2u^{2}-1}
\]
所以只需考察
\(\varphi(u)=2u-\sqrt{2u^{2}-1}\)，\(u\geqslant\frac1{\sqrt2}\).
当 \(u>\frac1{\sqrt2}\) 时，
\[
\varphi'(u)=2-\frac{2u}{\sqrt{2u^{2}-1}}
\]
当 \(\frac1{\sqrt2}<u<1\) 时，\(\sqrt{2u^{2}-1}<u\)，故 \(\varphi'(u)<0\)；当 \(u>1\) 时，\(\sqrt{2u^{2}-1}>u\)，故 \(\varphi'(u)>0\). 因此 \(\varphi\) 在 \(u=1\) 处取得最小值. 此时 \(|h|=|k|=1\)，且应取 \(h\)，\(k\) 同号，\(r^{2}=2\). 所求圆为
\[
(x-1)^{2}+(y-1)^{2}=2\quad\text{或}\quad (x+1)^{2}+(y+1)^{2}=2
\]
\end{solution}
